\chapter{Заключение}

В работе были рассмотрены спектральные характеристики внутренних представлений трёх семейств моделей компьютерного зрения: ResNet, ViT и Swin. Основной интерес состоял не в анализе итогового класса, предсказываемого моделью, а в изменении частотной структуры признаков внутри модели. Цель состояла в том, чтобы определить, по-разному ли архитектуры перераспределяют энергию между низкими, средними и более высокими пространственными частотами.

Для этого был собран единый протокол спектрального анализа. В ResNet анализировались карты признаков основных стадий. В ViT последовательность патчевых токенов переводилась обратно в двумерную решётку, при этом CLS-токен исключался, так как он не соответствует отдельной пространственной позиции. В Swin учитывалась стадийная структура модели и изменение размера токенной решётки между уровнями. После получения представлений строился спектр мощности, затем выполнялся переход к радиальному описанию. Для графиков использовалась средняя мощность в радиальном кольце, а для численных метрик --- суммарная энергия кольца. Это различие важно: именно суммарная энергия позволяет корректно считать долю низкочастотной энергии, поскольку радиальные кольца содержат разное число точек.

Отдельно была уточнена методика сравнения частотных радиусов. Частотная ось приводилась к общему масштабу, связанному с начальным слоем ResNet, а расчёт выполнялся внутри круга Найквиста. Из-за этого на части графиков у правой границы появляется небольшой подъём хвоста. В работе он не интерпретируется как содержательное усиление высоких частот. Это является визуальным артефактом радиального представления и нормировки кривых по максимуму. Численные выводы при этом опираются не на край графика, а на центроид, долю энергии в низкочастотной области и коэффициент экспоненциального сжатия.

Первое рабочее предположение относилось к свёрточным сетям. Ожидалось, что ResNet по мере перехода к глубоким стадиям будет смещать спектр представлений к меньшим частотным радиусам. Это предположение в целом подтвердилось. Во всех рассмотренных ResNet спектральный центроид снижается от начальных уровней к глубоким: у ResNet18 --- с $0.431\pm0.034$ до $0.055\pm0.002$, у ResNet34 --- с $0.436\pm0.036$ до $0.056\pm0.002$, у ResNet50 --- с $0.480\pm0.034$ до $0.071\pm0.004$, у ResNet101 --- с $0.484\pm0.032$ до $0.067\pm0.004$. Доля низкочастотной энергии также растёт: на глубоких уровнях она достигает $100\%$ при выбранном пороге $\rho<0.15$. Это значение не означает, что вся содержательная структура стала «низкочастотной» в обычном смысле. Оно связано ещё и с тем, что после приведения частотной оси к масштабу \texttt{layer1} доступный радиальный диапазон глубоких карт целиком попадает ниже выбранного порога.

Коэффициент экспоненциального спектрального сжатия у ResNet оказался самым высоким среди трёх семейств: от $1.930 \pm 0.109$ до $2.028 \pm 0.085$. Иными словами, по таблицам ResNet действительно демонстрирует наиболее устойчивое снижение центроида по глубине. Но здесь есть важное уточнение. Контрольный эксперимент показал, что сужение спектра нельзя полностью свести к обычному уменьшению пространственного разрешения. Реальный \texttt{layer4} ResNet18 оказался заметно уже, чем \texttt{layer1}, искусственно приведённый к размеру \texttt{layer4}. Следовательно, в ResNet есть эффект сверх простого понижения разрешения. При этом внутристадийный контроль ResNet50 показал, что внутри \texttt{layer3}, где пространственное разрешение фиксировано, кривые меняются гораздо спокойнее. Поэтому корректная формулировка такая: низкочастотная динамика ResNet возникает как совместный результат архитектурной иерархии и изменения самих признаков.

Второе предположение касалось ViT. Здесь ожидалось, что трансформер не обязан повторять поведение CNN, потому что он работает с последовательностью патчевых токенов и сохраняет одну и ту же патчевую решётку по глубине. Результаты это подтвердили. У ViT нет устойчивого экспоненциального сжатия спектра. У Tiny центроид почти возвращается к начальному значению: $0.120\pm0.011$ в начале и $0.121\pm0.004$ на глубоком блоке. У Small и Base максимум центроида приходится на средний уровень. У Large центроид растёт с $0.121\pm0.007$ до $0.154\pm0.007$.

Доля низкочастотной энергии у ViT также не даёт единой монотонной картины. У Tiny она остаётся примерно на одном уровне, у Small и Base снижается на среднем уровне и затем частично восстанавливается, а у Large падает с $62.201\pm4.225\%$ до $39.819\pm4.989\%$. Коэффициент $\lambda$ у ViT близок к нулю или отрицателен: от $-0.103 \pm 0.132$ у Tiny до $-0.334 \pm 0.078$ у Large. По этим числам ViT нельзя описать как архитектуру с последовательным низкочастотным сжатием. Более точный вывод звучит иначе: при фиксированной патчевой решётке ViT немонотонно перераспределяет спектральную энергию между блоками.

Важно не делать из этого лишних причинных выводов. Более широкий хвост у поздних блоков ViT не доказывает сам по себе, что модель «лучше сохраняет полезные детали». Графики показывают форму спектрального распределения, а не полезность конкретных частот для решения задачи. Возможные объяснения через механизм внимания, остаточные связи или MLP-блоки выглядят правдоподобно, но в этой работе они не проверялись отдельными абляционными экспериментами. Поэтому в выводе фиксируется наблюдение, а не причинная схема.

Третье предположение относилось к Swin. Предполагалось, что из-за иерархической структуры и операций объединения патчей Swin будет ближе к CNN на глубоких стадиях, но при этом его промежуточная динамика может отличаться от ResNet. Эта гипотеза подтвердилась с оговоркой. У Swin центроид к глубоким блокам заметно снижается: финальные значения лежат в диапазоне $0.072$--$0.078$. Доля низкочастотной энергии на глубоких уровнях также достигает $100\%$ при выбранном пороге. Коэффициент $\lambda$ положителен и находится в диапазоне от $0.951 \pm 0.353$ до $1.774 \pm 0.426$. По этим числам Swin ближе к ResNet, чем к ViT.

При этом контроль понижения разрешения уточняет вывод. В Swin значительная часть финального низкочастотного профиля объясняется уменьшением токенной решётки через \texttt{слияние патчей}. Реальная кривая и искусственно уменьшенное начальное представление не совпадают полностью, но основной низкочастотный пик у них близок. Поэтому нельзя утверждать, что всё финальное сжатие Swin вызвано только изменением признаков внутри трансформерных блоков. Таким образом, Swin занимает промежуточное положение: на финальных уровнях он приходит к низкочастотному профилю, но часть этого эффекта связана с архитектурным уменьшением разрешения.

Совместный анализ графиков и таблиц показывает, что три семейства моделей различаются не только числом параметров или способом обучения. Они действительно по-разному меняют частотный состав внутренних представлений. ResNet даёт наиболее последовательное снижение центроида. ViT показывает немонотонную динамику на фиксированной решётке патчей. Swin приходит к низкочастотному финальному профилю, но этот результат нужно читать вместе с контролем на уменьшение разрешения.

Практическая значимость работы состоит в том, что предложенный протокол можно использовать как инструмент диагностики внутренних представлений. Он помогает увидеть, насколько быстро модель концентрирует энергию в низкочастотной области, сохраняет ли среднечастотный хвост и где именно по глубине происходят наиболее сильные изменения. Для задач, где пространственная структура признаков особенно значима, такой анализ может быть полезен как дополнительный способ сравнения архитектур. Речь может идти, например, о медицинских изображениях, спутниковых снимках, дефектоскопии, текстурном анализе и задачах с мелкими объектами.

У работы есть ограничения. Спектральная энергия не равна информации в строгом смысле. Она не показывает напрямую, полезен ли конкретный частотный диапазон для классификации. В работе также не анализируется фаза спектра и не проверяется причинная связь между частотным профилем и итоговой точностью модели. Кроме того, в работе не проводились статистические тесты значимости различий между архитектурными семействами: значения среднее $\pm$ стандартное отклонение характеризуют разброс по изображениям, но не доказывают статистическую значимость различий. Ещё одно ограничение связано с ViT и Swin: перевод токенов в двумерную решётку нужен для сопоставления с CNN, но такая решётка не является полной копией свёрточной карты признаков. Токен уже содержит агрегированное описание патча, и это нужно учитывать при интерпретации.

Дальнейшее развитие работы возможно в нескольких направлениях. Во-первых, целесообразно исследовать, как спектральные метрики меняются не только у готовых моделей, но и в процессе обучения. Во-вторых, можно отдельно проверить причинные объяснения: например, сравнить варианты ViT и Swin с изменёнными механизмами внимания, оконной обработкой и объединением патчей. В-третьих, имеет смысл перенести анализ на детекцию и сегментацию, где пространственная локализация признаков имеет большее значение, чем в обычной классификации. Ещё одно направление --- анализ связи спектрального профиля с устойчивостью к шуму и сдвигам распределения данных.

Главный итог работы можно сформулировать так: архитектура задаёт не только конечный спектральный профиль представления, но и сам путь перераспределения энергии по глубине. ResNet демонстрирует наиболее устойчивое низкочастотное сжатие, ViT --- немонотонное изменение спектрального состава при фиксированной патчевой решётке, Swin --- финальный низкочастотный профиль, который частично объясняется уменьшением пространственного разрешения. Именно это различие и является основным результатом проведённого исследования.
